\section{A conditional complete-contour margin}\label{sec:contour}
This appendix consolidates the conditional result of RGD Note08, whose
normalization and scope were checked in E04 r4. The canonical scalar
hypotheses below remain unproved here. No new positivity domain follows.

For $x>0$, let $\mu=xC'(x)/C(x)$ and
$v=\mathrm d\mu/\mathrm d\log x$. They are the mean and variance of
$c_jx^j/C(x)$. Positivity and the infinite coefficient support imply
$v>0$, $\mu(0+)=0$ and $\mu(x)\to\infty$; thus there is a unique
$x=x(k)$ with $\mu(x)=k$ for each integer $k\ge1$. Set
$\phi(\theta)=e^{-ik\theta}C(xe^{i\theta})/C(x)$.
Coefficient extraction and determinant integration give
\begin{equation}\label{eq:contour-int}
 D_{n,k}=\frac{C(x)^nx^{-nk}}{n!(2\pi)^n}
 \int_{[-\pi,\pi]^n}\prod_{i=1}^n\phi(\theta_i)
 \prod_{i<j}4\sin^2\frac{\theta_i-\theta_j}{2}\dd\boldsymbol\theta.
\end{equation}
This standard identity also appears in Katkova's Lemma~1
\cite{Katkova0505174v1}; the integral is real but not pointwise positive.

\begin{proposition}[Conditional joint-order allowance]\label{prop:contour}
Fix $0<c\le1$, $B>0$ and $0<\delta_0\le1$. At the saddle assume
\begin{align}
 \log\phi(\theta)&=-v\theta^2/2+R(\theta),&
 |R(\theta)|&\le Bv|\theta|^3\quad(|\theta|\le\delta_0),\label{eq:contour-H1}\\
 |\phi(\theta)|&\le e^{-cv\theta^2/2}&&(|\theta|\le\pi),\label{eq:contour-H2}
\end{align}
where the local logarithm is continuous and zero at the origin.
Put $A=8+2\log(2/c)$ and
\[
 G(n)=\max\left\{\frac{An^2}{c\delta_0^2},\,
 \frac{4An^4}{3c},\,16B^2(A/c)^3n^8\right\}.
\]
If $v\ge G(n)$, then
\begin{equation}\label{eq:contour-margin}
 D_{n,k}>\frac{C(x)^nx^{-nk}}{4n!(2\pi)^n}Z_n(v)>0,
 \qquad Z_n(v)=(2\pi)^{n/2}v^{-n^2/2}\prod_{j=1}^n j!.
\end{equation}
\end{proposition}
\begin{proof}
The Gaussian Vandermonde integral over $\R^n$ equals $Z_n(v)$.
Determinant integration and monic Hermite norms
$\sqrt{2\pi}\,j!v^{-j-1/2}$ prove this normalization, including the
$n!$ factor. Choose $\delta^2=An^2/(cv)$ and
$\tau=nBv\delta^3$. The assumptions give
$\delta\le\delta_0$, $\tau\le1/4$ and
$n(n-1)\delta^2/3\le1/4$. On $|\theta_i|\le\delta$, the real part
of the complete phase product is at least
$e^{-\tau}\cos\tau\,e^{-v\sum\theta_i^2/2}$.
The circular Vandermonde is at least the ordinary squared Vandermonde
times $S=(\sin\delta/\delta)^{n(n-1)}\ge e^{-1/4}$.
Here $\sin y/y\ge e^{-y^2/3}$ for $|y|\le1$, and all coincident
limits are continuous.

The Gaussian mass outside this cube is at most $TZ_n(v)$, where
$T=2^{n^2/2}e^{-v\delta^2/4}$. This follows by splitting the
Gaussian exponent in half on $\sum\theta_i^2\ge\delta^2$.
Using \eqref{eq:contour-H2}, the absolute contribution of the
\emph{entire} complementary contour region is at most $UZ_n(v)$, with
$U=(2/c)^{n^2/2}e^{-cv\delta^2/4}$, by the same argument at
variance parameter $cv/2$. This includes every mixed region; no
separate combinatorial exterior payment is omitted.
Our choice of $A$ gives $U=e^{-2n^2}$ and $T\le e^{-2n^2}$.
The real integral in \eqref{eq:contour-int}, divided by $Z_n(v)$, is
therefore at least
\[
 e^{-\tau}\cos\tau\,S(1-T)-U
 \ge e^{-1/2}\cos(1/4)(1-e^{-2})-e^{-2}
 >\frac{61}{224}>\frac14.
\]
This proves the claim.
\end{proof}

If canonical analysis were to establish the hypotheses for all $x\ge x_0$
with $\mu(x)/4\le v(x)\le\mu(x)$, then
$k\ge\lceil\max\{k_0,4G(n)\}\rceil$, for an integer
$k_0\ge\mu(x_0)$, would give an explicit sufficient domain.
Its $n^8$ scale is the cost of this allowance, not an optimal law.
The statement neither joins the complementary region nor constitutes
an induction at fixed $k$. The source literature already reports
stronger tail coverage; see the qualifications in the introduction.
